%%%%%%%%%%  scan idx 185  |  folio 172  %%%%%%%%%%
puisque $U'$ est affine et que l'on a $n-1\ge1$, d'où une contradiction.

\textbf{Remarque 2.2.6.1.} Sans l'hypothèse de séparation sur $X$, (2.2.6) pour un ouvert affine $U$ de $X$ serait faux, comme le montre l'exemple du plan affine où l'origine a été dédoublée (EGA I 5.5.11).

\textbf{Proposition 2.2.7.} Soit $S$ un schéma noethérien, séparé, localement factoriel (i.e.\ dont les anneaux locaux sont factoriels). Alors $S$ est divisoriel (2.2.5).

En particulier (EGA IV 21.11.1) :

\textbf{Corollaire 2.2.7.1.} Tout schéma noethérien régulier séparé est divisoriel.

\emph{Preuve} de 2.2.7. On peut supposer $S$ intègre. Soit $(S_i)_{i\in I}$ un recouvrement fini de $S$ par des ouverts affines non vides. En vertu de (2.2.6), $S-S_i$ est purement de codimension $1$, donc, d'après (EGA IV 21.7.4), est le support d'un diviseur positif. En d'autres termes, pour chaque $i\in I$, il existe un faisceau inversible $L_i$ et une section $f_i\in\Gamma(S,L_i)$ telle que $S_i=S_{f_i}$. La famille $(L_i)$ est alors ample (2.2.4) d'après (2.2.3 (i$'$)). Ceci achève la démonstration.

\textbf{Remarque 2.2.7.2.} Comme en (2.2.6), l'hypothèse de séparation sur $S$ est essentielle. En effet, prenons encore une fois pour $S$ le plan affine où l'origine a été dédoublée. Rappelons que $S$ est obtenu en recollant deux exemplaires du plan affine $\mathbb E^2_k$ ($k$, un corps) par l'identité sur le complémentaire $U$ de l'origine $O$ ; on a donc une projection canonique
\[
p:S\longrightarrow\mathbb E^2_k,
\]
qui induit un isomorphisme au-dessus de $U$, $p^{-1}(O)$ se composant de deux points $O_1$ et $O_2$. Un faisceau inversible sur $S$ s'obtient en recollant deux faisceaux inversibles sur $\mathbb E^2_k$ par un isomorphisme au-dessus de $U$ ; comme tout faisceau inversible sur $\mathbb E^2_k$ est trivial, un faisceau inversible sur $S$ est donc défini par un élément $g\in\Gamma(U,\mathcal O_U^*)$ ; mais $g$ est nécessairement une constante, puisque la restriction $\Gamma(\mathbb E^2_k,\mathcal O_{\mathbb E^2_k})\to\Gamma(U,\mathcal O_U)$ est un isomorphisme, $O$ étant de profondeur $2$ dans le plan ; donc tout faisceau
